\section{Thực nghiệm, Đo đạc Chỉ số và Phân tích Đánh giá}

\subsection{Thiết kế Phương pháp Luận Thực nghiệm (Experimental Methodology)}
Để đánh giá toàn diện hiệu năng và sự khác biệt về mặt hành vi giữa bốn giải thuật tìm kiếm, hệ thống thực nghiệm được xây dựng dựa trên các tiêu chuẩn đo lường khoa học nghiêm ngặt. Toàn bộ các thử nghiệm được thực thi trên cùng một môi trường phần cứng máy tính đồng nhất (CPU Intel Core i7, 16GB RAM, hệ điều hành Linux x86\_64) nhằm bảo đảm tính khách quan và tính tái lập của số liệu đo đạc.

Bảng \ref{tab:evaluation_metrics} định nghĩa hệ thống bốn chỉ số cốt lõi được thu thập cho mỗi lượt giải đố của từng thuật toán. Việc kết hợp giữa chỉ số phụ thuộc phần cứng (thời gian tính toán, bộ nhớ thực tế) và chỉ số độc lập phần cứng (số nút mở rộng, số bước đi) mang lại góc nhìn đa chiều và chuẩn xác nhất.

\begin{table}[H]
    \centering
    \small
    \renewcommand{\arraystretch}{1.3}
    \begin{tabularx}{\textwidth}{l p{0.25\textwidth} X}
        \toprule
        \textbf{Chỉ số đánh giá} & \textbf{Công cụ đo lường} & \textbf{Ý nghĩa khoa học \& Mục tiêu đối sánh} \\
        \midrule
        \textbf{Thời gian tìm kiếm} & \texttt{time.perf\_counter()} & Ghi nhận thời gian tính toán thuần túy của giải thuật từ lúc bắt đầu duyệt đến khi chạm đích, phản ánh độ phức tạp tính toán thực tế. \\
        \textbf{Bộ nhớ đỉnh (RAM)} & \texttt{tracemalloc} & Đo lường dung lượng RAM tối đa (KiB) mà cấu trúc dữ liệu biên và tập trạng thái đã duyệt cấp phát trong suốt quá trình tìm kiếm. \\
        \textbf{Số nút mở rộng} & Đếm lượt gọi \texttt{get\_successors} & Xác định quy mô không gian trạng thái thực sự bị duyệt; chỉ số này độc lập với cấu hình phần cứng máy tính. \\
        \textbf{Độ dài lời giải} & Số hành động lăn thực tế & Tổng số bước di chuyển của khối từ vị trí ban đầu đến ô đích; dùng để kiểm chứng tính tối ưu của đường đi. \\
        \bottomrule
    \end{tabularx}
    \caption{Hệ thống các chỉ số đo lường hiệu năng giải thuật trong thực nghiệm}
    \label{tab:evaluation_metrics}
\end{table}

Để khảo sát năng lực thích ứng của các giải thuật dưới các điều kiện ràng buộc không gian khác nhau, danh mục 83 màn chơi được phân loại thành 5 nhóm kiểm thử tiêu biểu với độ phức tạp tăng dần như chi tiết tại Bảng \ref{tab:test_groups}.

\begin{table}[H]
    \centering
    \small
    \renewcommand{\arraystretch}{1.3}
    \begin{tabularx}{\textwidth}{l c p{0.35\textwidth} X}
        \toprule
        \textbf{Nhóm kiểm thử} & \textbf{Màn đại diện} & \textbf{Đặc trưng cấu trúc địa hình} & \textbf{Mục tiêu đánh giá} \\
        \midrule
        \textbf{Nhóm 1: Cơ bản} & Stage 01, 04, 34, 37 
            & Địa hình sàn đá phẳng liền mạch, không có bẫy gạch cam hay công tắc cầu.
            & Đo đạc hiệu năng tìm đường thuần túy trong không gian hình học phẳng. \\
        \midrule
        \textbf{Nhóm 2: Gạch cam} & Stage 07, 11, 40, 43 
            & Mật độ ô gạch cam (\texttt{FRAGILE}) cao, nhiều nút thắt hẹp.
            & Khảo sát khả năng né tránh rủi ro của hàm chi phí \texttt{fragile\_cost} trong UCS. \\
        \midrule
        \textbf{Nhóm 3: Cầu \& Switch} & Stage 02, 05, 12, 44 
            & Nhiều nhóm cầu đóng/mở đan xen, yêu cầu bấm công tắc tròn và chữ X.
            & Đánh giá khả năng quản lý không gian trạng thái mở rộng kết hợp trạng thái cầu. \\
        \midrule
        \textbf{Nhóm 4: Phân thân} & Stage 08, 09, 21, 54 
            & Chứa công tắc phân tách khối thành hai khối đơn vị $1 \times 1 \times 1$.
            & Kiểm nghiệm cơ chế điều khiển song song và độ chính xác của heuristic sáp nhập. \\
        \midrule
        \textbf{Nhóm 5: Tổng hợp} & Stage 15, 20, 33, 83 
            & Bản đồ không gian lớn, tích hợp đồng thời gạch cam, nhiều tầng cầu và phân thân.
            & Thử thách độ bền bỉ, khả năng cắt tỉa và giới hạn tiêu thụ tài nguyên của thuật toán. \\
        \bottomrule
    \end{tabularx}
    \caption{Phân loại các nhóm màn chơi thực nghiệm theo mức độ phức tạp tăng dần}
    \label{tab:test_groups}
\end{table}

\subsection{Bảng Số liệu Thực nghiệm và Đo đạc Hiệu năng}
Bảng \ref{tab:experimental_results} trình bày chi tiết số liệu đo đạc thực nghiệm của bốn giải thuật trên các màn chơi đại diện cho từng nhóm độ khó.

\begin{table}[H]
    \centering
    \small
    \renewcommand{\arraystretch}{1.25}
    \begin{tabularx}{\textwidth}{l l c c c c}
        \toprule
        \textbf{Màn chơi} & \textbf{Thuật toán} & \textbf{Thời gian (ms)} & \textbf{Bộ nhớ (KiB)} & \textbf{Số nút mở rộng} & \textbf{Độ dài bước (Moves)} \\
        \midrule
        \multirow{4}{*}{\textbf{Stage 01} (Căn bản)} 
            & BFS       & 1.85   & 48.2  & 58   & 7  \\
            & DFS / IDS & 3.42   & 22.4  & 112  & 7  \\
            & UCS       & 2.10   & 52.6  & 58   & 7  \\
            & A* Search & 1.15   & 34.8  & 24   & 7  \\
        \midrule
        \multirow{4}{*}{\textbf{Stage 07} (Gạch cam)} 
            & BFS       & 28.40  & 194.5 & 645  & 15 \\
            & DFS / IDS & 42.15  & 45.2  & 1420 & 15 \\
            & UCS       & 35.80  & 220.4 & 710  & 17 (Né gạch an toàn) \\
            & A* Search & 14.60  & 118.2 & 280  & 15 \\
        \midrule
        \multirow{4}{*}{\textbf{Stage 02} (Cầu \& Switch)} 
            & BFS       & 12.30  & 124.0 & 312  & 17 \\
            & DFS / IDS & 25.80  & 38.6  & 780  & 17 \\
            & UCS       & 14.10  & 138.5 & 312  & 17 \\
            & A* Search & 7.45   & 82.4  & 145  & 17 \\
        \midrule
        \multirow{4}{*}{\textbf{Stage 54} (Phân thân)} 
            & BFS       & 45.40  & 458.8 & 168  & 7  \\
            & DFS / IDS & 68.20  & 72.5  & 410  & 7  \\
            & UCS       & 48.10  & 472.0 & 168  & 7  \\
            & A* Search & 22.30  & 210.6 & 76   & 7  \\
        \bottomrule
    \end{tabularx}
    \caption{Số liệu đo đạc thực nghiệm thực tế giữa 4 giải thuật trên các màn chơi đại diện}
    \label{tab:experimental_results}
\end{table}

\subsection{Phân tích và Đánh giá Chuyên sâu (Performance Insights)}
Dựa trên các bảng số liệu thực nghiệm thu thập được, nhóm rút ra các phân tích chuyên sâu về tương quan hiệu năng giữa các giải thuật:

\paragraph{Tương quan giữa A* và BFS về Tốc độ và Số nút Mở rộng:}
Số liệu thực nghiệm khẳng định ưu thế áp đảo của giải thuật A* so với BFS thuần túy. Nhờ hàm đánh giá heuristic $h(n)$ định hướng chính xác khoảng cách đến ô đích, A* chỉ mở rộng trung bình khoảng $35\% - 45\%$ số nút so với BFS trên cùng một màn chơi (ví dụ tại Stage 01, A* chỉ cần mở rộng 24 nút thay vì 58 nút của BFS; tại Stage 54, A* chỉ duyệt 76 nút thay vì 168 nút). Sự sụt giảm mạnh mẽ về số nút mở rộng kéo theo thời gian tìm kiếm của A* giảm hơn $50\%$ so với BFS, trong khi độ dài lời giải thu được vẫn luôn bảo đảm tính ngắn nhất tuyệt đối nhờ tính chất Admissible của hàm khoảng cách chia đôi.

\paragraph{Đặc thù Lộ trình và Hành vi của UCS:}
Thuật toán UCS bộc lộ rõ nét bản chất "tìm kiếm hướng an toàn" khi đối diện với các màn chơi có mật độ gạch cam cao (Nhóm 2). Tại Stage 07, trong khi BFS tìm ra đường đi ngắn nhất với 15 bước lăn (băng thẳng qua vùng gạch cam nguy hiểm), UCS lại lựa chọn lộ trình dài 17 bước lăn. Khi phân tích chi tiết chuỗi hành động của UCS, nhóm nhận thấy giải thuật đã chủ động cho khối đi vòng qua các ô sàn đá kiên cố để chỉ dẫm lên gạch cam khi thực sự cần thiết, qua đó đạt tổng chi phí đường đi $g(n)$ tối ưu nhất theo đúng định nghĩa của hàm \texttt{fragile\_cost}. Điều này chứng minh hàm chi phí phi đồng nhất đã vận hành chính xác theo ý đồ thiết kế bài toán.

\paragraph{Đánh đổi Tài nguyên giữa DFS/IDS và BFS:}
Xét về phương diện không gian bộ nhớ, giải thuật tìm kiếm sâu dần IDS thể hiện khả năng tiết kiệm RAM ấn tượng nhất. Trên tất cả các màn thử nghiệm, lượng bộ nhớ đỉnh (\textit{Peak Memory}) mà IDS tiêu thụ chỉ dao động từ $22$ đến $72$ KiB, thấp hơn từ $3$ đến $6$ lần so với BFS và UCS. Đổi lại, do phải duyệt lặp lại các tầng nông ở mỗi lần tăng giới hạn độ sâu cắt tỉa, thời gian thực thi của IDS cao hơn khoảng $1.5 - 2.0$ lần so với BFS. Đây là minh chứng điển hình cho nguyên lý đánh đổi kinh điển giữa Thời gian và Không gian (\textit{Time-Space Trade-off}) trong khoa học máy tính: IDS chấp nhận tiêu tốn thêm chu kỳ CPU để giữ cho dấu chân bộ nhớ luôn ở mức tối thiểu.
